\documentclass{article}

\usepackage[utf8]{inputenc}

\usepackage{biblatex}
\addbibresource{main.bib}

\begin{document}

\title{Synthèse bibliographique}

\section{Éléments de définition du geste technique}

Afin de débuter ma documentation sur le thème de l'apprentissage des gestes techniques, je me suis appuyé sur le livre de Lucile Vadcard \cite{vadcard_gestes_2022}.
L'ouvrage permet en particulier de dégager des éléments de définition des gestes techniques en tant qu'objet scientifique.
\newline

Plusieurs définitions sont proposées au long de l'ouvrage. Dans son sens le plus global, le geste est défini comme un \emph{"ensemble coordonné de mouvements du corps, participant à une action finalisée"}. Cette définition permet de différencier le geste technique du "geste amical" ainsi qu'une simple exécution physique, mais limite le geste à son étude quantitative, approche que Roux et Bril (2002) ont montré limitée pour saisir les déterminants. Leplat (2013) fournit un exemple concret de l'échec de cette approche en formation technique (soudure).

Il s'agit donc de considérer le geste comme un \emph{"moyen produit par le sujet pour atteindre un objectif, compte tenu des contraintes de l'environnement"}: le geste va donc prendre en compte à la fois le sujet (personne exécutant), la tâche donnée ainsi que l'environnement auquel le geste devra s'adapter (Bril et Goasdoué, 2009).
\newline

Le geste technique peut être étudié sous plusieurs caractères:
\begin{itemize}
\item Son caractère \emph{moteur}, le plus souvent représenté repose sur la coordination des systèmes de posture et de mouvement. L'approche ergonomique désigne que Berthoz, dans son ouvrage de 1997, permet de faire évoluer ce caractère vers une approche située du geste et non seulement comme un mouvement en montrant que le sujet simule cognitivement afin de prévoir les événements à venir à partir de ses connaissances. L'approche anthropologique du geste technique, menée notamment par Bril, désigne Paillard (1976) comme un fondateur de ce placement du geste en situation. Le caractère moteur est particulièrement étudié par les sciences du vivant.
\item Son caractère \emph{fonctionnel} place l'accent sur la finalité du geste. Ainsi, un même mouvement peut avoir des fonctions différentes (e.g. biologique, geste-outil, communication, socioculturelle). Une autre catégorisation peut être proposée entre \emph{geste-action} (destiné à produire un effet) et \emph{geste-signe} (destiné à communiquer), ou entre \emph{invariance} et \emph{adaptabilité} rappelant les travaux de Vergnaud (1996) sur les schèmes.
\item Son caractère \emph{structurant} fait le lien avec les connaissances du sujet, en rappelant la notion de \emph{cerveau simulateur} de Berthoz. Ce caractère rejoint notamment l'approche de Piaget (1967) distinguant les connaissances tirées des objets eux-mêmes et celles tirées des actions du sujet sur les objets, ainsi que celle des schèmes de Vergnaud. Il permet notamment des adaptations vis-à-vis de l'environnement.
\end{itemize}
Le caractère structurant est rappelé par Mellet d'Huart et Michel \cite{mellet-dhuart_faciliter_2005} comme primordial dans l'apprentissage car le but de l'apprentissage est d'engendrer de nouvelles connaissances à partir de la formation reçue à partir des connaissances préalables du sujet.

Plusieurs dimensions se dégagent du geste technique et permet d'étayer sa définition:
\begin{itemize}
    \item La \emph{planification}: le geste est réfléchi, divisé en étapes et parfois coordonné au sein d'une équipe.
    \item La \emph{responsabilité}: le sujet doit pouvoir juger son geste sur sa validité, tant pour le geste en lui-même que pour l'ensemble du processus dans lequel s'intègre le geste.
    \item L'\emph{histoire sociale}: le geste est effectué au sein d'un contexte construit par le passé.
    \item L'\emph{histoire personnelle}: Le geste est appris selon le vécu des sujets, les répétitions, l'appropriation du geste.
    \item L'\emph{instrumentation}: le geste technique est médié par un objet technique (instrument), ou le corps lui-même.
\end{itemize}

L'instrumentation donne lieu à la distinction entre deux fonctions de l'objet technique: l'\emph{instrument} (permet de percevoir) et l'\emph{outil} (permet d'agir). Rabardel introduit la notion de "genèse instrumentale" pour établir les relations entre objet technique et sujet: l'\emph{instrumentation} désigne l'action de l'objet technique sur le sujet par le biais de l'affordance (possibilité offertes) et réciproquement, l'\emph{instrumentalisation}, qui correspond à la sélection des possibilités par le sujet (voire l'ajout de fonctionnalité secondaires). La \emph{genèse instrumentale} correspond à la cohérence créée par le sujet entre l'objet et l'activité.
On se rend compte de l'interaction entre les deux par le système perception-action-conceptualisation \cite{mellet-dhuart_faciliter_2005}, et leur relations avec le caractère structurant du geste.

La sophistication de l'instrumentation donne également lieu à une progressive \emph{dématérialisation} des situations de travail et donc de nouvelles formes de matérialité plus cognitives que motrices, nécessitant un modèle mental permettant d'anticiper l'action de la machine pour pouvoir la paramétrer.
\section{Approches de l'apprentissage}

Plusieurs domaines se sont intéressés à l'apprentissage des gestes, donnant lieu à des approches distinctes, bien que non mutuellement exclusives et ne devant pas être considérées en comparaison les unes des autres, mais en les combinant:
\begin{itemize}
    \item L'approche \emph{physiologique}, fondée à partir des travaux de Marey et Bernstein notamment, se focalise sur le caractère moteur du geste et accorde un rôle prédominant à la \emph{répétition}. Elle mesure pour évaluer le geste l'écart référent/référé en prenant un geste idéal comme référent.
    \item L'approche \emph{anthropologique}, à partir des travaux de Leroi-Gourhan et Mauss, ne s'intéresse pas au caractère moteur du geste mais cherche plutôt à le replacer dans un contexte socioculturel. Elle utilise plutôt l'apprentissage par l'\emph{imitation}.
    \item L'approche \emph{ergonomique} recherche les marges de manoeuvre nécessaires pour préserver la santé du sujet et accorde une place de choix à la \emph{simulation} dans la formation.
    \item L'approche \emph{didactique} se focalise sur le caractère structurant à partir des cadres théoriques de Vergnaud (schèmes mettant en relation les buts/anticipations, règles d'action, possibilité d'inférences en situation, moments opératoires ; boucle de rétroaction structurante au sein d'une situation didactique) et du constructivisme développé par Piaget (assimilation: répétition et généralisation ; accomodation: adaptation quand l'assimilation est mise à mal)
\end{itemize}
Des combinaisons peuvent apparaître, par exemple sur le caractère moteur en faisant se suivre un processus imitatif (amorce) puis répétitif (perfectionnement).

Mellet d'Huart et Michel \cite{mellet-dhuart_faciliter_2005} rappellent l'importance de considérer la situation d'apprentissage comme complexe avec des facteurs relevant de l'apprenant, du contenu, du formateur, de l'institution (dont l'importance est soulignée par Petit et Oudart \cite{petit_apprendre_2017} car influençant les choix didactiques), de l'environnement.
Pour Pastré, cité par eux, la dynamique des compétences à viser est de passer d'une \emph{coordination agie} (geste) à sa \emph{coordination conceptualisée}, en passant par des processus de \emph{désincorporation} et de \emph{décontextualisation}.

Deux voies complémentaires sont abordées pour la formation: le \emph{tutorat} et les \emph{simulations} avec leurs problématiques et limites respectives.
\newline

Dans l'approche du tutorat (e.g. compagnonnage, alternance, médecine), des questions se posent déjà concernant l'encadrement du tutoré: quel cadre légal et quelles libertés autorisées en pratique? Comment concilier le tutorat avec les contraintes de productivité du travail.

Le tutorat se base sur le \emph{couple maître-apprenti}, bien qu'entouré par le collectif et que ce couple soit rendu instable par les multiples cycles de stage. Ce couple permet en particulier, par l'autorité conférée au maître via son excellence une \emph{imitation prestigieuse} (Mauss). Mais le tutorat ne se limite pas à cette dimension proche du caractère moteur: il s'inscrit également dans un dimension d'histoire collective et personnelle en permettant à l'apprenti de se voir en position de maîtrise au sein d'un corps de métier et d'un groupe social. Le collectif notamment apporte des connaissances connexes participant au caractère structurant du geste: des règles de conduits, des qualités personnelles attendues.

Les temporalités sont également importante dans le cadre d'un tutorat: le temps long permet la construction du couple maître/apprenti mais aussi de la coutume de travail comme l'explique Balacheff (1988), et l'apprentissage initial du geste par la démonstration expliquée et les actions guidées. Le temps court, quant à lui, permet des négociations de l'action selon le contexte, participant à l'adaptabilité et donc au caractères fonctionnels et structurants.
De plus, un point de vue épistémologique met l'accent sur la nécessaire éducation à l'attention (quand être attentif): la transmission n'est donc pas une simple somme d'occasions d'apprentissage.

Cette temporalité concerne également la progressivité des tâches confiées à l'apprenti. Cette progressivité dépend notamment de l'expertise du maître, de la nature des gestes.
Kunegel (2007) fournit un \emph{modèle de la fonction tutorale} on donne pour composante les formes de l'interaction du tutelle (sélection des tâches, évaluation) et les configuration tutorales (degré d'assistance de mise en travail)

En réalité, on se rend aussi compte que l'aspect moteur n'est pas toujours abordé ou que dans certaines temporalités (boucles de feedback courtes).
Actuellement les approches pour le tutorat    sont plutôt dirigées vers la neurophysiologie, cognitivisme ou béhaviorisme, en opposition avec une approche constructiviste laissant plus de place à la didactique.
\newline

L'approche par simulation est actuellement très étudiée (voir Granry, Gallagher, O'Sullivan), notamment en lien avec des philosophies de formation prônées par la HAS ("jamais la première fois sur le patient") et des développement technologiques offrant de plus en plus de possibilités en termes de \emph{réalisme}. La \emph{simulation} est définie comme un \emph{"ensemble de techniques pour la formation pratique, lorsqu'elle ne se déroule pas au travail"}.

En réalité, cette approche prônant la supériorité du réalisme déplace la réflexion pédagogique vers le formateur qui devra adapter son scénario pédagogique. De plus, le peu d'études s'intéressant à la comparaison avec d'autres méthodes ne permettent pas de conclure à une supériorité de l'approche par simulation.

Il est à rappeler également que la simulation n'implique pas nécessairement un objet technique et que, s'il y a objet technique, il ne s'agit pas nécessairement d'un simulateur, qui est défini comme \emph{"la partie tangible de la simulation qui  expose le processus simulé}.
Le simulateur cherche à reproduire les lois comportementales et peut altérer le réalisme à des fins de simplifications permettant la construction d'une meilleure situation d'apprentissage (patients standardisés, soudure).

Dans une simulation, différents éléments sont à mettre en correspondance avec les caractères du geste: les élements de contexte (fonctionnel), les éléments tangibles (moteur) et les connaissances du sujet (structurant). La caractère dynamique de la situation permet ensuite de restructurer les connaissances du sujet.

Afin d'étudier les caractéristiques des simulations, Balacheff (1994) introduit le domaine de validité épistémologique: quelles informations initiales sont données au sujet ? Quelles actions lui sont autorisées ? Quels effets ont ces actions ? Quelles variations sont possibles ?
On peut alors établir une typologie des objets techniques de simulation selon leurs similitudes: fortuites (animaux, humains, cadavres), statiques (maquettes) et dynamique (qui sont dès lors des simulateurs: jeux de rôles, maquettes motorisées...).

Pastré propose notamment, en se focalisant sur le couple sujet/simulation, de prendre des problèmes extraits de situations de travail pour en faire des situations d'apprentissage.

Dans un cadre se voulant plus didactique, les EIAH (environnements informatiques pour l'apprentissage humain) se focalise sur l'intention pédagogique.
La simulation, dans ce cadre, fait l'objet de critiques concernant la place qui occupent les apprenants: ceux-là ne peuvent comprendre le rapport référence/représentant, les enseignants se retrouvent donc à compenser à leur place, donc approche moins constructiviste.

La notion de \emph{fidélité} (physique, fonctionnelle, psychologique) plutôt que de réalisme permet d'évoluer en ce sens, et l'apport de l'épistémologie permet de supposer un critère de pertinence d'une simulation: les connaissances qu'elle mobilise sont-elles bien les mêmes que dans la situation de travail ?
\newline

Enfin, l'approche didactique, fondée notamment sur l'importance de l'erreur, reporte le travail sur la création de situations d'apprentissage adaptées. Pastré semble être un auteur fondamental dans cette approche.
Mellet d'Huart et Michel \cite{mellet-dhuart_faciliter_2005} soulignent par exemple l'importance de comprendre les fonctionnements et dysfonctionnements cognitifs des apprenants durant la formation, en plus de l'analyse experte de la tâche, du travail sur le scénario pédagogique.

Ils proposent une approche utilisant le \emph{réalité virtuelle} afin de faire se rejoindre la simulation et l'approche didactique. Cette possibilité offre en particulier la possibilité d'une démarche constructiviste comme souligé par Roussos, Johnon, Moher et al. (1999). Winn (2003) soulignait également les avantages permis par l'affordance de la réalité virtuelle.

La VR permet en particulier un guidage, une séparation des paramètres sensitivo-moteurs, un feedback plus direct, plus de répétitions. Cela permet de construire au préalable un coordination agie pour ensuite se focaliser, en réel, sur la conceptualisation.

\section{Approches de l'évaluation}

Nous cherchons ici à étudier comment les moyens d'évaluation prennent en considération les différents caractères du geste technique (moteur, fonctionnel, structurant) dans différents domaines.
\newline

Tout d'abord, afin d'apporter des éléments théoriques sur l'évaluation, nous pouvons citer une typologie proposée par De Ketele \cite{de_ketele_ne_2010}. Dans cet article publié principalement à destination du milieu scolaire, l'auteur rappelle les 3 objectifs possibles de l'évaluation:
\begin{itemize}
    \item l'\emph{orientation} d'une nouvelle action
    \item la \emph{régulation}, amélioration, d'une action en courtes
    \item la \emph{certification} d'une action passée, toujours sociale et donc engageante du point de vue de l'évaluateur.
\end{itemize}
La plupart du temps, l'évaluation dans le cadre de l'apprentissage est certificative, ce qui peut apporter des problématiques sur la place donnée à l'\emph{erreur}. De Ketele donne deux valeurs possibles à cette dernière: elle peut être \emph{régulatrice} et donc moteur de l'apprentissage, comme le suggère l'approche didactique des apprentissages expliquée plus haut.
Elle peut au contraire être associée à une valeur \emph{certificative} et sanctionne alors.

Au-delà de son objectif, l'évaluation en pratique peut se décliner sous diverses modalités:
\begin{itemize}
    \item \emph{sommative}: il s'agit d'une approche par note, souvent demandée. Celle-ci peut être attribuée par addition ou soustraction
    \item \emph{descriptive}: en identifiant et décrivant les caractéristiques de l'objet à évaluer, l'évaluateur fournit un feedback particulièrement intéressant dans des objectifs de régulation ou d'orientation. Si cette approche est rejetée dans l'éducation, elle est toutefois plus appréciée en milieu professionnel, donc dans le domaine visé par ce mémoire.
    \item \emph{herméneutique} (ou autres dénominations selon les cadres théoriques), basée sur un recueil d'indices variés, puis leur organisation en vue de valider une hypothèse parmi d'autres. C'est le cas dans beaucoup de situations nécessitant une orientation de décision, ou alors dans le cadre professionnalisant, des stages.
\end{itemize}
Historiquement, une confusion a été faite par Scriven (1967) entre finalité et modalité de l'évaluation en comparant sommatif et formatif (régulatif), associant ainsi sommatif et certificatif, ce qui peut encore aujourd'hui biaiser certains cadres évaluatifs.
\newline

Pour De Ketele, le \emph{processus évaluatif} consiste à confronter un \emph{référent} et un \emph{référé} afin de fonder une prise de décision. Quoique d'autres auteurs voie un jugement comme finalité de l'évaluation, De Ketele place ce jugement porté comme une partie du processus et non sa finalité. En effet, les étapes du processus évaluatif, à savoir 1) le recueil d'informations, 2) leur confrontation à des critères préétablis, 3) l'attribution d'une signification à ce résultat, 4) la prise de décision fondée ; révèlent un système de valeurs de l'évaluateur dans le choix des informations recensées et des critères retenus.

L'éablissement du processus évaluatif doit garder en tête plusieurs questions:
\begin{itemize}
    \item Pertinence des informations recueillies
    \item Validité du processus: mesure-t-on ce que l'on déclare mesurer ?
    \item Fiabilité (invariabilité des résultats en fonction de l'évaluateur)
\end{itemize}

Enfin, De Ketele offre une discussion sur la distinction entre \emph{ressources} et \emph{compétences} et leurs lien: la maîtrise d'une compétence mobilise des ressources, mais l'évaluation peut porter sur les ressources (qui sont alors explicitement demandées ou indiquées par le guidage) ou sur les compétences, pluridimmensionnelles, où l'évalué sera laissé libre des ressources qu'il mobilise. En cela la compétence est plus grande que la somme des ressources qu'elle mobilise (le tout est différent de la somme des parties).
De là, il propose une procédure pour évaluer les compétences, passant notamment par l'identification de critères minimaux et de perfectionnement que l'on peut relier, dans le cadre d'une évaluation des gestes, à leurs caractères \emph{fonctionnels} et \emph{sructurants}.

En médecine, la thèse de Camille Faudeux \cite{faudeux_sfp_2014} \cite{faudeux_formation_2015} propose une validation de plusieurs grilles d'évaluation disponibles en annexe. On constate que ces grilles présentent seulement des éléments de mouvements (caractère moteur).
En restant dans le domaine médical, le référentiel technique fourni aux équipes pédagogiques PSC \cite{direction_generale_de_la_securite_civile_et_de_la_gestion_des_risques_references_2026} inclut également des éléments de justification portant plutôt sur le caractère fonctionnel des gestes.
Dans la même philosophie, les recommandations ERC de 2025 \cite{greif_european_2025} placent les gestes de la réanimation cardiaque dans le cadre plus général de la prise en charge globale du patient.
\newline

Le règlement de la licence STAPS de l'Université Paris-Est Créteil \cite{universite_paris-est_creteil_livret_2022} ainsi que les grilles d'évaluation des APSA en lycée \cite{ministere_de_leducation_nationale_referentiel_nodate} peuvent faire intervenir des critères à la fois moteurs et fonctionnels, selon les disciplines (les pratiques d'expression en STAPS sont évaluées sur des caractères moteurs tandis que le handball par exemple l'ait sur des critères plutôt fonctionnels)
\newline

L'évaluation du caractère structurant est toutefois peu abordé dans ces référentiels d'évaluation.

\printbibliography

\end{document}